\section{PPO：能用，但不轻}

\subsection{从 policy gradient 到 PPO}

课上先快速回顾了 PPO 的理论谱系：
\begin{enumerate}
    \item 直接 policy gradient：概念最简单，但方差很高。
    \item TRPO：用重要性采样和 trust region 控制 policy 漂移。
    \item PPO：把更新约束成 clip 形式，既稳定又容易实现。
\end{enumerate}

核心目标可以写成标准的 clipped objective：
\[
\mathcal{L}_{\mathrm{PPO}}(\theta)
=
\mathbb{E}\left[
\min\Big(r_t(\theta) A_t,\ \mathrm{clip}(r_t(\theta),1-\epsilon,1+\epsilon)A_t\Big)
\right].
\]

\begin{figure}[H]
\centering
\includegraphics[width=0.95\textwidth]{slides-images/slide-017.jpg}
\caption{PPO 的理论回顾：policy gradient、TRPO、PPO clip}
\end{figure}

\begin{figure}[H]
\centering
\includegraphics[width=0.95\textwidth]{slides-images/slide-019.jpg}
\caption{PPO 在概念上其实不难：就是一个带裁剪的目标函数}
\end{figure}

\subsection{工程上为什么麻烦}

课堂里强调，PPO 的“理论版本”和“代码版本”差距很大。真正麻烦的地方主要有三类：
\begin{itemize}
    \item \textbf{rollouts} 很贵。每次更新都要跑模型采样，生成样本是最慢的那部分。
    \item \textbf{value function} 让系统更重。它要额外占显存，还引入了训练调参压力。
    \item \textbf{reward shaping} 和 \textbf{GAE} 让实现细节变多。虽然都不是神秘概念，但凑在一起就很容易把训练管线搞复杂。
\end{itemize}

\begin{figure}[H]
\centering
\includegraphics[width=0.95\textwidth]{slides-images/slide-021.jpg}
\caption{语言模型里的 PPO 近似成一个 bandit：token 动作序列，末端一个 dense reward}
\end{figure}

\begin{figure}[H]
\centering
\includegraphics[width=0.95\textwidth]{slides-images/slide-023.jpg}
\caption{PPO 的一个实际实现：outer loop 先采样，再在 inner loop 里做梯度更新}
\end{figure}

\begin{figure}[H]
\centering
\includegraphics[width=0.95\textwidth]{slides-images/slide-024.jpg}
\caption{PPO 的 loss 计算：policy loss、value loss、clip range 等部分要一起维护}
\end{figure}

\subsection{reward shaping 和 GAE}

在语言模型里，PPO 往往被当成一个 contextual bandit：输入一个 prompt，输出一个整段回答，最后再看 reward。为了让学习更稳定，实践里经常把 reward 分解成两类：
\begin{itemize}
    \item per-token 的 KL 惩罚，用来约束新 policy 不要离 reference 太远。
    \item 末端的任务成功 reward，用来表示真正的目标是否完成。
\end{itemize}

而 GAE 的作用，是把高方差的原始回报，变成一个更平滑的 advantage 估计。课堂里的关键点不是公式本身，而是它的工程意义：\textbf{如果没有这些技巧，PPO 很容易训练得又慢又不稳。}

\begin{figure}[H]
\centering
\includegraphics[width=0.95\textwidth]{slides-images/slide-026.jpg}
\caption{reward shaping：per-token KL 惩罚加上末端任务 reward}
\end{figure}

\begin{figure}[H]
\centering
\includegraphics[width=0.95\textwidth]{slides-images/slide-027.jpg}
\caption{GAE：用 advantage 去替代直接的 reward，降低梯度方差}
\end{figure}

